%%%%%%%%%%%%%%%%%%%%%%%%%%%%%%%%%%%%%%%%%%%%%%%%%%%%%%%%%%%%%%%%%%%%%%%%%%%%%%%
\section{Scientific Questions}
\label{sec:scientific_questions}

The thesis is guided by the following main question:

\begin{quote}
\textbf{Can the transport properties of organic thin-film transistors be measured
in a way that describes the accumulated semiconductor sheet, rather than being
dominated by contact effects?}
\end{quote}

This question is divided into three subquestions:
\begin{description}
    \item[\textbf{Q1 -- Field-effect operation:}] Do the fabricated devices and
    the measurement setup show stable field-effect operation with sufficiently
    low gate leakage?
    \item[\textbf{Q2 -- Contact influence:}] How strongly do contact resistance
    and charge injection affect conventional transfer and output characteristics,
    including threshold-voltage estimates?
    \item[\textbf{Q3 -- Four-terminal method and dynamic response:}] How doese the gated van
    der Pauw method give a  contact-independent sheet conductance,
    mobility as well as background
    resistance, and do transient or circuit effects have to be considered?
\end{description}

The first part of the work tests the experimental platform. The prototype setup
must provide stable current driving and reliable voltage detection, while the
fabricated devices must show low gate leakage, reproducible field-effect
modulation, and ordered transfer and output curves. Without these conditions,
neither conventional transistor analysis nor gated van der Pauw analysis can give
meaningful transport parameters.

The second part examines contact influence. In two-terminal measurements, the
same electrodes inject current and define the measured voltage, so injection
barriers and contact resistance can change the slope and shape of the curves. The
gated van der Pauw measurement tests whether a sheet conductance obtained from
separate voltage probes is more robust. The resistance of the full
current-injection path is analysed separately to compare the gate-dependent
channel response with any gate-independent background contribution.

The third part addresses dynamic effects. Localized states in organic
semiconductors can trap and release carriers over long time scales, affecting
threshold-voltage extraction, hysteresis, and circuit response
\cite{Hirwa2015TrapStates,Bobbert2012OperationalStability}. DLTS-type transient
measurements and a simple amplifier test are therefore used to check whether
slow relaxation and external circuit elements influence the device response.

Three working hypotheses follow from these questions. First, in a suitable
accumulation regime, the gated van der Pauw sheet conductance should increase
approximately linearly with effective gate bias, and the extracted mobility
should be similar for different sufficiently small measurement currents. Second,
two-terminal analysis should be more sensitive to contact and series effects than
the four-terminal sheet-conductance extraction. Third, transient and amplifier
measurements should reveal whether slow relaxation or external circuit time
constants must be considered when interpreting nominally static measurements.
